% Part: first-order-logic
% Chapter: natural-deduction
% Section: proving-things

\documentclass[../../../include/open-logic-section]{subfiles}

\begin{document}

\iftag{FOL}
      {\olfileid[fr]{fol}{ntd}{pro}}
      {\olfileid[fr]{pl}{ntd}{pro}}

\olsection{Exemples de \usetoken{p}{derivation}}

\begin{ex}
Construisons une !!{derivation} de l'!!{sentence} $(!A \land !B) \lif !A$.

Commençons par écrire la conclusion recherchée tout en bas de la
!!{derivation}.
\begin{prooftree}
\AxiomC{}
\UnaryInfC{$(!A\land !B) \lif !A$}
\end{prooftree}

Il faut ensuite déterminer quel type d'inférence pourrait avoir
pour conclusion un !!{sentence} de cette forme. L'!!{main operator}
de la conclusion est $\lif$; essayons donc d'y parvenir par la règle
\Intro{\lif}. Il est préférable d'écrire les hypothèses utilisées
et d'étiqueter les règles d'inférence au fur et à mesure : on verra
ainsi facilement si toutes les hypothèses ont été !!{discharged}s
à la fin de la preuve.
\begin{prooftree}
\AxiomC{$\Discharge{!A \land !B}{1}$}
\DeduceC{$!A$}
\DischargeRule{\Intro{\lif}}{1}
\UnaryInfC{$(!A\land !B) \lif !A$}
\end{prooftree}

Il reste à compléter les étapes qui conduisent de l'hypothèse
$!A \land !B$ à $!A$. Le seul connecteur à traiter étant $\land$,
nous utilisons sa règle d'élimination. On obtient la preuve suivante :
\begin{prooftree}
\AxiomC{$\Discharge{!A \land !B}{1}$}
\RightLabel{\Elim{\land}}
\UnaryInfC{$!A$}
\DischargeRule{\Intro{\lif}}{1}
\UnaryInfC{$(!A\land !B) \lif !A$}
\end{prooftree}
Nous avons maintenant une !!{derivation} correcte de
$(!A \land !B) \lif !A$.
\end{ex}

\begin{ex}
Construisons maintenant une !!{derivation} de
$(\lnot !A \lor !B) \lif (!A \lif !B)$.

Commençons par écrire la conclusion recherchée tout en bas de la
!!{derivation}.
\begin{prooftree}
\AxiomC{}
\UnaryInfC{$(\lnot !A \lor !B) \lif (!A \lif !B)$}
\end{prooftree}
Pour trouver une règle logique qui puisse donner cette conclusion,
examinons les connecteurs qui y figurent : $\lnot$, $\lor$ et
$\lif$. Seule la première occurrence de $\lif$ nous intéresse pour
l'instant, car elle est l'!!{main operator} de l'!!{sentence}
final. Les connecteurs $\lnot$, $\lor$ et la seconde occurrence de
$\lif$ sont dans la portée d'un autre connecteur; nous les traiterons
donc plus tard. Commençons par la règle \Intro{\lif}. Une
application correcte doit avoir la forme suivante :
\begin{prooftree}
\AxiomC{$\Discharge{\lnot !A \lor !B}{1}$}
\DeduceC{$!A \lif !B$}
\DischargeRule{\Intro{\lif}}{1}
\UnaryInfC{$(\lnot !A \lor !B) \lif (!A \lif !B)$}
\end{prooftree}
Deux possibilités s'offrent alors à nous. Nous pouvons poursuivre
de bas en haut en cherchant une autre application de \Intro{\lif},
ou travailler de haut en bas en appliquant \Elim{\lor}. Choisissons
la seconde possibilité. L'hypothèse $\lnot !A \lor !B$ sera la
prémisse la plus à gauche de \Elim{\lor}. Pour que l'application
soit correcte, les deux autres prémisses doivent être identiques
à la conclusion $!A \lif !B$. Chacune peut cependant être dérivée
à partir d'une hypothèse supplémentaire : l'un des deux membres de
la disjonction $\lnot !A \lor !B$. Notre !!{derivation} aura donc
la forme suivante :
\begin{prooftree}
\AxiomC{$\Discharge{\lnot !A \lor !B}{1}$}
\AxiomC{$\Discharge{\lnot !A}{2}$}
\DeduceC{$!A \lif !B$}
\AxiomC{$\Discharge{!B}{2}$}
\DeduceC{$!A \lif !B$}
\DischargeRule{\Elim{\lor}}{2}
\TrinaryInfC{$!A \lif !B$}
\DischargeRule{\Intro{\lif}}{1}
\UnaryInfC{$(\lnot !A \lor !B) \lif (!A \lif !B)$}
\end{prooftree}

Dans chacune des deux branches de droite, nous voulons
!!{derive} $!A \lif !B$. La règle \Intro{\lif} est tout indiquée :
\begin{prooftree}
\AxiomC{$\Discharge{\lnot !A \lor !B}{1}$}
\AxiomC{$\Discharge{\lnot !A}{2}, \Discharge{!A}{3}$}
\DeduceC{$!B$}
\DischargeRule{\Intro{\lif}}{3}
\UnaryInfC{$!A \lif !B$}
\AxiomC{$\Discharge{!B}{2}, \Discharge{!A}{4}$}
\DeduceC{$!B$}
\DischargeRule{\Intro{\lif}}{4}
\UnaryInfC{$!A \lif !B$}
\DischargeRule{\Elim{\lor}}{2}
\TrinaryInfC{$!A \lif !B$}
\DischargeRule{\Intro{\lif}}{1}
\UnaryInfC{$(\lnot !A \lor !B) \lif (!A \lif !B)$}
\end{prooftree}

Pour compléter les deux parties manquantes, il nous faut dériver
$!B$ à partir de $\lnot !A$ et $!A$ au milieu, et à partir de
$!A$ et $!B$ à droite. Commençons par la branche du milieu :
$\lnot !A$ et $!A$ sont les deux prémisses de \Elim{\lnot}.
\begin{prooftree}
\AxiomC{$\Discharge{\lnot !A}{2}$}
\AxiomC{$\Discharge{!A}{3}$}
\RightLabel{\Elim{\lnot}}
\BinaryInfC{$\lfalse$}
\DeduceC{$!B$}
\end{prooftree}
La règle \FalseInt{} permet alors d'obtenir la conclusion $!B$
et d'achever cette branche.
\begin{prooftree}
\AxiomC{$\Discharge{\lnot !A \lor !B}{1}$}
\AxiomC{$\Discharge{\lnot !A}{2}$}
\AxiomC{$\Discharge{!A}{3}$}
\RightLabel{\Elim{\lnot}}
\BinaryInfC{$\lfalse$}
\RightLabel{\FalseInt}
\UnaryInfC{$!B$}
\DischargeRule{\Intro{\lif}}{3}
\UnaryInfC{$!A \lif !B$}
\AxiomC{$\Discharge{!B}{2}, \Discharge{!A}{4}$}
\DeduceC{$!B$}
\DischargeRule{\Intro{\lif}}{4}
\UnaryInfC{$!A \lif !B$}
\DischargeRule{\Elim{\lor}}{2}
\TrinaryInfC{$!A \lif !B$}
\DischargeRule{\Intro{\lif}}{1}
\UnaryInfC{$(\lnot !A \lor !B) \lif (!A \lif !B)$}
\end{prooftree}

Passons à la branche la plus à droite. Il faut ici se rappeler que
la définition d'une !!{derivation} \emph{permet de décharger des
hypothèses}, sans \emph{l'imposer}. Autrement dit, il est permis
de dériver $!B$ à partir d'une seule des hypothèses $!A$ et $!B$,
sans utiliser l'autre. Or dériver $!B$ à partir de~$!B$ est
immédiat : $!B$ seul constitue une telle !!{derivation}, sans
aucune inférence. Nous pouvons donc supprimer l'hypothèse~$!A$.
\begin{prooftree}
\AxiomC{$\Discharge{\lnot !A \lor !B}{1}$}
\AxiomC{$\Discharge{\lnot !A}{2}$}
\AxiomC{$\Discharge{!A}{3}$}
\RightLabel{\Elim{\lnot}}
\BinaryInfC{$\lfalse$}
\RightLabel{\FalseInt}
\UnaryInfC{$!B$}
\DischargeRule{\Intro{\lif}}{3}
\UnaryInfC{$!A \lif !B$}
\AxiomC{$\Discharge{!B}{2}$}
\RightLabel{\Intro{\lif}}
\UnaryInfC{$!A \lif !B$}
\DischargeRule{\Elim{\lor}}{2}
\TrinaryInfC{$!A \lif !B$}
\DischargeRule{\Intro{\lif}}{1}
\UnaryInfC{$(\lnot !A \lor !B) \lif (!A \lif !B)$}
\end{prooftree}
Dans la !!{derivation} achevée, l'inférence \Intro{\lif} la plus
à droite ne décharge donc aucune hypothèse.
\end{ex}

\begin{ex}
Jusqu'ici, nous n'avons pas eu besoin de la règle \FalseCl{}.
Sa particularité est de permettre la décharge d'une hypothèse
qui n'est pas une sous-!!{formula} de la conclusion de la règle.
Elle est étroitement liée à \FalseInt{} : cette dernière en est
un cas particulier. Il existe une logique, appelée « logique
intuitionniste », dans laquelle seule \FalseInt{} est admise parmi
ces deux règles. On recourt à \FalseCl{} en dernier ressort,
lorsque les autres tentatives échouent. Supposons par exemple
que nous voulions !!{derive} $!A \lor \lnot !A$. Notre stratégie
habituelle consisterait à essayer de !!{derive} $!A \lor \lnot !A$
par $\Intro{\lor}$. Il faudrait alors !!{derive} soit $!A$, soit
$\lnot !A$, sans aucune hypothèse, ce qui n'est pas possible
en général. La règle \FalseCl{} permet de sortir de cette impasse :
\begin{prooftree}
  \AxiomC{$\Discharge{\lnot(!A \lor \lnot !A)}{1}$}
  \DeduceC{$\lfalse$}
  \DischargeRule{\FalseCl}{1}
  \UnaryInfC{$!A \lor \lnot !A$}
\end{prooftree}
Nous cherchons maintenant une !!{derivation} de $\lfalse$ à partir
de $\lnot(!A \lor \lnot !A)$. Puisque $\lfalse$ est la conclusion
de $\Elim{\lnot}$, essayons cette règle :
\begin{prooftree}
  \AxiomC{$\Discharge{\lnot(!A \lor \lnot !A)}{1}$}
  \DeduceC{$\lnot !A$}
  \AxiomC{$\Discharge{\lnot(!A \lor \lnot !A)}{1}$}
  \DeduceC{$!A$}
  \RightLabel{\Elim{\lnot}}
  \BinaryInfC{$\lfalse$}
  \DischargeRule{\FalseCl}{1}
  \UnaryInfC{$!A \lor \lnot !A$}
\end{prooftree}
Pour trouver une !!{derivation} de~$\lnot !A$, notre stratégie
nous conduit à appliquer~$\Intro{\lnot}$ :
\begin{prooftree}
  \AxiomC{$\Discharge{\lnot(!A \lor \lnot !A)}{1}, \Discharge{!A}{2}$}
  \DeduceC{$\lfalse$}
  \DischargeRule{\Intro{\lnot}}{2}
  \UnaryInfC{$\lnot !A$}
  \AxiomC{$\Discharge{\lnot(!A \lor \lnot !A)}{1}$}
  \DeduceC{$!A$}
  \RightLabel{\Elim{\lnot}}
  \BinaryInfC{$\lfalse$}
  \DischargeRule{\FalseCl}{1}
  \UnaryInfC{$!A \lor \lnot !A$}
\end{prooftree}
On obtient ici facilement $\lfalse$ en appliquant $\Elim{\lnot}$
à l'hypothèse $\lnot(!A \lor \lnot !A)$ et à $!A \lor \lnot !A$,
qui résulte de notre nouvelle hypothèse $!A$ par~$\Intro{\lor}$ :
\begin{prooftree}
  \AxiomC{$\Discharge{\lnot(!A \lor \lnot !A)}{1}$}
  \AxiomC{$\Discharge{!A}{2}$}
  \RightLabel{\Intro{\lor}}
  \UnaryInfC{$!A \lor \lnot !A$}
  \RightLabel{\Elim{\lnot}}
  \BinaryInfC{$\lfalse$}
  \DischargeRule{\Intro{\lnot}}{2}
  \UnaryInfC{$\lnot !A$}
  \AxiomC{$\Discharge{\lnot(!A \lor \lnot !A)}{1}$}
  \DeduceC{$!A$}
  \RightLabel{\Elim{\lnot}}
  \BinaryInfC{$\lfalse$}
  \DischargeRule{\FalseCl}{1}
  \UnaryInfC{$!A \lor \lnot !A$}
\end{prooftree}
À droite, on suit la même stratégie, mais on obtient $!A$ par~\FalseCl :
\begin{prooftree}
  \AxiomC{$\Discharge{\lnot(!A \lor \lnot !A)}{1}$}
  \AxiomC{$\Discharge{!A}{2}$}
  \RightLabel{\Intro{\lor}}
  \UnaryInfC{$!A \lor \lnot !A$}
  \RightLabel{\Elim{\lnot}}
  \BinaryInfC{$\lfalse$}
  \DischargeRule{\Intro{\lnot}}{2}
  \UnaryInfC{$\lnot !A$}
  \AxiomC{$\Discharge{\lnot(!A \lor \lnot !A)}{1}$}
  \AxiomC{$\Discharge{\lnot !A}{3}$}
  \RightLabel{\Intro{\lor}}
  \UnaryInfC{$!A \lor \lnot !A$}
  \RightLabel{\Elim{\lnot}}
  \BinaryInfC{$\lfalse$}
  \DischargeRule{\FalseCl}{3}
  \UnaryInfC{$!A$}
  \RightLabel{\Elim{\lnot}}
  \BinaryInfC{$\lfalse$}
  \DischargeRule{\FalseCl}{1}
  \UnaryInfC{$!A \lor \lnot !A$}
\end{prooftree}
\end{ex}

\begin{prob}
Donner des !!{derivation}s établissant les résultats suivants :
\begin{enumerate}
\item $!A \land (!B \land !C) \Proves (!A \land !B) \land !C$.
\item $!A \lor (!B \lor !C) \Proves (!A \lor !B) \lor !C$.
\item $!A \lif (!B \lif !C) \Proves !B \lif (!A \lif !C)$.
\item $!A \Proves \lnot\lnot !A$.
\end{enumerate}
\end{prob}

\begin{prob}
Donner des !!{derivation}s établissant les résultats suivants :
\begin{enumerate}
\item $(!A \lor !B) \lif !C \Proves !A \lif !C$.
\item $(!A \lif !C) \land (!B \lif !C) \Proves (!A \lor !B) \lif !C$.
\item $\Proves \lnot(!A \land \lnot !A)$.
\item $!B \lif !A \Proves \lnot !A \lif \lnot !B$.
\item $\Proves (!A \lif \lnot !A) \lif \lnot !A$.
\item $\Proves \lnot(!A \lif !B) \lif \lnot !B$.
\item $!A \lif !C \Proves \lnot (!A \land \lnot !C)$.
\item $!A \land \lnot !C \Proves \lnot (!A \lif !C)$.
\item $!A \lor !B, \lnot !B \Proves !A$.
\item $\lnot !A \lor \lnot !B \Proves \lnot(!A \land !B)$.
\item $\Proves (\lnot !A \land \lnot !B) \lif\lnot(!A \lor !B)$.
\item $\Proves \lnot(!A \lor !B) \lif (\lnot !A \land \lnot !B)$.
\end{enumerate}
\end{prob}

\begin{prob}
Donner des !!{derivation}s établissant les résultats suivants :
\begin{enumerate}
\item $\lnot(!A \lif !B) \Proves !A$.
\item $\lnot(!A \land !B) \Proves \lnot !A \lor \lnot !B$.
\item $!A \lif !B \Proves \lnot !A \lor !B$.
\item $\Proves \lnot \lnot !A \lif !A$.
\item $!A \lif !B, \lnot !A \lif !B \Proves !B$.
\item $(!A \land !B) \lif !C \Proves (!A \lif !C) \lor (!B \lif !C)$.
\item $(!A \lif !B) \lif !A \Proves !A$.
\item $\Proves (!A \lif !B) \lor (!B \lif !C)$.
\end{enumerate}
(Tous ces résultats nécessitent la règle~$\FalseCl$.)
\end{prob}

\begin{editorial}
Dans le deuxième exemple, la source parle improprement de séquent
final et situe à gauche la branche issue de $!A$ et $!B$, qui est
à droite. Dans l'arbre intermédiaire suivant l'emploi de \FalseInt,
l'étiquette \Intro{\lfalse} a été corrigée en \Elim{\lnot},
conformément aux prémisses et à l'explication. Enfin, l'impossibilité
de dériver $!A$ ou $\lnot !A$ sans hypothèse est une affirmation
générale sur le schéma : certaines instances particulières sont
démontrables. Les exercices restent des schémas pour des énoncés
arbitraires.
\end{editorial}

\end{document}
